% ============================================================
% 30_layout_spacing.tex — Layout-Grundlagen, Spacing-Skala, Box-Bindungen
% ============================================================

\setlength{\parindent}{0pt}

\widowpenalty=10000
\clubpenalty=10000
\brokenpenalty=4000
\predisplaypenalty=10000

\raggedcolumns
\newcounter{zsfFirstChap}\setcounter{zsfFirstChap}{1}

%
%
%
%
%
%
%
%
%
%
%
%
%
%
%
%
%
%
%
\providecommand{\ZSFDensityFactor}{1}

% ── \ZSFderive: eine Ableitung, die man wiederholen kann ─────────────
% Setzt das Register sofort UND merkt sich die Rechnung. \ZSFChapterDensity
% (unten) laesst die gemerkte Kette in derselben Reihenfolge noch einmal
% laufen und bekommt damit den ganzen Registersatz auf einen neuen Faktor.
%
% Der Grund fuer die Kette statt einer Handvoll Neuberechnungen: Ein Register
% ist eine LAENGE, kein Ausdruck. \ZSFboxPadX hat den Wert, den \ZSFspaceS bei
% seiner Berechnung hatte — die vier Basisregister nachtraeglich zu aendern
% erreicht es nicht mehr. Wer eine neue abgeleitete Groesse anlegt, benutzt
% deshalb \ZSFderive und nicht \setlength; sie faehrt dann von selbst mit.
\makeatletter
\newcommand{\ZSF@spacingChain}{}
\newcommand{\ZSFderive}[2]{%
  \setlength{#1}{#2}%
  \expandafter\gdef\expandafter\ZSF@spacingChain\expandafter{%
    \ZSF@spacingChain\setlength{#1}{#2}}%
}
% Der Dokumentwert, auf den \ZSFChapterDensityReset zurueckstellt. Muss VOR der
% ersten Ableitung gesichert werden, sonst ist er beim Zuruecksetzen schon weg.
\let\ZSF@densityDocument\ZSFDensityFactor

%
%
%
%
%
%
%
%
%
%
%
%
\newcommand{\ZSFChapterDensity}[1]{\def\ZSFDensityFactor{#1}\ZSF@spacingChain
  \ZSFApplyDisplaySkips}
\newcommand{\ZSFChapterDensityReset}{%
  \let\ZSFDensityFactor\ZSF@densityDocument \ZSF@spacingChain
  \ZSFApplyDisplaySkips}
\makeatother

% ── Display-Skips: die EINE Stelle, und sie muss wiederholbar sein ───
% Die engen Abstände um Formeln sind eine ZSF-Entscheidung. LaTeX legt sie
% aber in den Groessenbefehlen ab: \normalsize und \footnotesize setzen alle
% vier Skips auf die KLASSENWERTE zurueck. Jeder Schriftwechsel im Inhalt
% loescht damit die Entscheidung — lautlos, denn Formeln sehen auch mit den
% Klassenwerten korrekt aus, nur luftiger.
%
% Deshalb ist die Politik ein benanntes Makro und kein einmaliges \setlength:
% Wer die Groesse wechselt, ruft sie danach erneut auf. Genutzt von
% \ZSFfontContent/\ZSFfontContentDense (der `font`-Regler), vom
% \AtBeginDocument in 70_document_settings und von \ZSFChapterDensity oben.
% \jot und \arraycolsep laufen hier mit, obwohl kein Groessenbefehl sie
% zuruecksetzt: Sie gehoeren zur selben Frage — »wie viel Luft hat eine
% Formel« — und eine zweite Anwendungsstelle waere eine zweite Auslegung.
\newcommand{\ZSFApplyDisplaySkips}{%
  \setlength{\abovedisplayskip}{\ZSFspaceXS}%
  \setlength{\belowdisplayskip}{\ZSFspaceXS}%
  \setlength{\abovedisplayshortskip}{\ZSFspaceXS}%
  \setlength{\belowdisplayshortskip}{\ZSFspaceXS}%
  \setlength{\mathsurround}{\ZSFspaceXS}%
  \setlength{\jot}{\ZSFmathRowSep}%
  \setlength{\arraycolsep}{\ZSFmathColSep}%
}

\newlength{\ZSFspaceXS}\ZSFderive{\ZSFspaceXS}{\ZSFDensityFactor\dimexpr 1pt\relax}
\newlength{\ZSFspaceS} \ZSFderive{\ZSFspaceS}{\ZSFDensityFactor\dimexpr 4pt\relax}
\newlength{\ZSFspaceM} \ZSFderive{\ZSFspaceM}{\ZSFDensityFactor\dimexpr 7pt\relax}
\newlength{\ZSFspaceL} \ZSFderive{\ZSFspaceL}{\ZSFDensityFactor\dimexpr 10pt\relax}

%
%
%
%
%
%
%
%
%
%
%
%
%
%
%
%
%
%
%
\providecommand{\ZSFDensityBoxes}{1}
\providecommand{\ZSFDensityText}{1}
\providecommand{\ZSFDensityTables}{1}
\providecommand{\ZSFDensityStructure}{1}

% Outer (Text-Box-Bindung) und Sep (Trennlinien) — beide folgen der Dichte.
% Outer ist bewusst als \ZSFspaceM abgeleitet statt erneut ausgeschrieben:
% es IST der Standard-Aussenabstand, kein eigener Wert.
\newlength{\ZSFspaceOuter}\ZSFderive{\ZSFspaceOuter}{\ZSFspaceM}
\newlength{\ZSFspaceSep}  \ZSFderive{\ZSFspaceSep}{\ZSFDensityBoxes\dimexpr
  \ZSFDensityFactor\dimexpr 3pt\relax\relax}

% Absatzabstand im Fliesstext — der grösste einzelne vertikale Posten einer
% ZSF und deshalb der wichtigste Teil der Dichte-Stellschraube. Muss nach den
% Basisregistern stehen, damit \ZSFDensityFactor bereits bekannt ist.
% Der Bereichsfaktor greift über das Zwischenregister: Ein Faktor vor einem
% Skip-Register skaliert Soll-, Streck- und Schrumpfanteil gemeinsam; zwei
% Faktoren direkt hintereinander kann TeX dagegen nicht lesen.
\newlength{\ZSFparskipBase}
\ZSFderive{\ZSFparskipBase}{\ZSFDensityFactor\dimexpr 3.5pt\relax
                    plus \ZSFDensityFactor\dimexpr 1pt\relax
                   minus \ZSFDensityFactor\dimexpr 0.5pt\relax}
\ZSFderive{\parskip}{\ZSFDensityText\ZSFparskipBase}

% ── Zeilenhöhe: die zweite globale Stellschraube ─────────────────────
% Bewusst NICHT an \ZSFDensityFactor gekoppelt. Abstände sind Leerraum und
% vertragen jede Skalierung; Zeilenhöhe enthält die Schrift selbst — sie mit
% demselben Faktor zu stauchen, lässt Formeln und Akzente kollidieren.
% Wer zusätzlich zur Dichte enger setzen will, dreht hier separat:
%   \def\ZSFLeadingFactor{0.92}   % in preamble.tex, vor diesem \input
\providecommand{\ZSFLeadingFactor}{0.95}
\linespread{\ZSFLeadingFactor}

% ------------------------------------------------------------
% Spalten-Layout
% ------------------------------------------------------------
% \columnsep folgt der Dichte bewusst NICHT: Beim Verdichten sollen die
% Zeilen enger stehen, nicht die Spalten zusammenrücken — ein schmalerer
% Steg macht die ZSF schlechter lesbar, obwohl er kaum Platz spart.
\newlength{\ZSFcolumnSep}\setlength{\ZSFcolumnSep}{10pt}
\setlength{\columnsep}{\ZSFcolumnSep}
\setlength{\multicolsep}{\ZSFspaceS}
\setlength{\emergencystretch}{2em}

% ------------------------------------------------------------
% Schwellwerte und Titelbalken-Hilfsgrössen
% ------------------------------------------------------------
% Mindesthöhe, die vor einem Titelbalken frei sein muss — sonst bricht die
% Spalte davor um. Die Reserve muss den Balken UND den Anfang der folgenden
% Box tragen, sonst bleibt der Balken allein am Spaltenende stehen.
% Deshalb staffeln sich die Werte nach der Grösse der eröffneten Einheit:
% Kapitel > Abschnitt > Unterabschnitt.
% Die drei Werte sind EINE Staffelung, kein Trio unabhängiger Zahlen. Wer die
% ZSF früher oder später umbrechen lassen will, dreht deshalb den Faktor und
% nicht die drei Register einzeln — nur so bleibt ihr Verhältnis gewahrt.
% Grösser = die Spalte bricht eher vorzeitig um und lässt
% Balken seltener knapp am Fuss stehen; kleiner = dichter gefüllte Spalten.
\providecommand{\ZSFBreakReserveFactor}{1}
\newlength{\ZSFbreakThresholdChapter}
  \setlength{\ZSFbreakThresholdChapter}{\ZSFBreakReserveFactor\dimexpr 10\baselineskip\relax}
\newlength{\ZSFbreakThreshold}
  \setlength{\ZSFbreakThreshold}{\ZSFBreakReserveFactor\dimexpr 7\baselineskip\relax}
\newlength{\ZSFbreakThresholdSubsubsection}
  \setlength{\ZSFbreakThresholdSubsubsection}{\ZSFBreakReserveFactor\dimexpr 5\baselineskip\relax}
\newlength{\ZSFtitleBarHeight}\setlength{\ZSFtitleBarHeight}{9.2pt}
\newlength{\ZSFsubsectionBarBeforeSkip}\ZSFderive{\ZSFsubsectionBarBeforeSkip}{\ZSFDensityStructure\ZSFspaceL}
\newlength{\ZSFsubsubsectionBarBeforeSkip}\ZSFderive{\ZSFsubsubsectionBarBeforeSkip}{\ZSFDensityStructure\ZSFspaceS}
\newlength{\ZSFsubsectionAfterChapterGap}\ZSFderive{\ZSFsubsectionAfterChapterGap}{\ZSFDensityStructure\ZSFspaceXS}

% ── Der Abstand NACH einem Titelbalken ───────────────────────────────
% EIN Register für alle vier Balken, und es ist der GANZE Abstand von der
% Unterkante des Balkens bis zur Oberkante des folgenden Blocks — Fliesstext
% wie Box. Der Weg dorthin ist \ZSFInterlude (unten), aufgerufen im
% after-Hook \ZSFTitleBarAfter (60_boxes).
%
% Vorher war dieser Abstand herrenlos, und zwar nach demselben Muster wie
% \ZSFBoxBeforeSkip (styles/README → »Abstand hat genau einen Schreiber«):
% Zwei Register (\ZSFchapterBarAfterSkip 2pt, \ZSFsubsectionAfterGap 1pt)
% setzten je EINEN Posten von dreien. Dazu kamen \parskip (3.5pt, weil der
% Text nach dem Balken ein neuer Absatz ist) und die Zeilenschaltung
% (\baselineskip − Höhe der ersten Zeile). Gemessen am Satz: Balken → Box
% 1pt, Balken → Fliesstext 6.8pt, Kapitelbalken → Fliesstext 9.5pt. Drei
% Abstände an derselben Stelle des Systems, keiner davon gewählt.
%
% Der Wert (3pt) liegt zwischen den beiden Zuständen von vorher: enger als der
% Fliesstext-Fall (6.8/9.5pt, den niemand gewählt hatte), etwas luftiger als
% der Box-Fall (1pt, der nur deshalb so knapp war, weil dort kein \parskip
% dazukam). Am Satz geprüft, nicht gerechnet: 2pt drückt den ersten Satz an
% den Balken, 4pt bringt die alte Lücke zurück.
%
% Die Staffelung Kapitel > Abschnitt ist bewusst aufgegeben: Die Ebene wird
% über Balkenhöhe, Farbe und den Abstand DAVOR unterschieden
% (\ZSFsubsectionBarBeforeSkip), nicht über einen Punkt danach.
\providecommand{\ZSFBarGapFactor}{1}
\newlength{\ZSFbarAfterGap}\ZSFderive{\ZSFbarAfterGap}{\ZSFBarGapFactor\dimexpr
  \ZSFDensityStructure\dimexpr\ZSFDensityFactor\dimexpr 3pt\relax\relax\relax}
% Der sichtbare Aussenradius ist Form, nicht Leerraum: Auch bei dichterem Satz
% bleiben die Kanten bewusst gleich weich. Der Innenradius wird aus der
% jeweiligen Rahmenbreite abgeleitet, damit die Kontur in der Kurve nicht
% dicker wird als auf den geraden Kanten.
\newlength{\ZSFboxArc}\setlength{\ZSFboxArc}{3pt}
\newlength{\ZSFboxRuleWidth}\setlength{\ZSFboxRuleWidth}{0.5pt}
\newlength{\ZSFboxInnerArc}\ZSFderive{\ZSFboxInnerArc}{\dimexpr
  \ZSFboxArc-\ZSFboxRuleWidth\relax}
\newlength{\ZSFgoalConditionRuleWidth}\setlength{\ZSFgoalConditionRuleWidth}{1.5pt}
\newlength{\ZSFgoalConditionBorderlineWidth}\setlength{\ZSFgoalConditionBorderlineWidth}{0.2pt}
\newlength{\ZSFgoalConditionInnerArc}\ZSFderive{\ZSFgoalConditionInnerArc}{\dimexpr
  \ZSFboxArc-\ZSFgoalConditionRuleWidth\relax}
% Abstand zwischen zwei Gliedern einer Zielkette (Bedingung, Schritt, Ziel).
% Ein eigenes Register, weil der Abstand in der goalbox eine ENTSCHEIDUNG ist
% und nicht das Nebenprodukt der Zeilenschaltung: Die Glieder sind verschieden
% hoch (Pille, Displayzeile, gerahmtes Ziel), und ohne einen benannten Wert
% ergibt sich der Rhythmus aus der Tinte statt aus dem Satz (siehe
% \ZSFInterlude). Gleich \ZSFboxPadY, damit der Abstand zum Rahmen derselbe
% ist wie der zwischen den Gliedern — die Kette liest sich dadurch als ein
% gleichmässiger Stapel.
\newlength{\ZSFgoalLinkSep}  \ZSFderive{\ZSFgoalLinkSep}{\ZSFDensityBoxes\ZSFspaceS}
% Zeilenhöhe der ZSFtable in drei Stufen — der Regler `rows` wählt zwischen
% ihnen (20_tables). Drei Werte statt eines, weil die Höhe einer Zeile nicht
% aus der Tabellenart folgt, sondern aus ihrem Inhalt: Eine Tabelle mit
% \dfrac-Zellen braucht mehr Höhe, ein Symbolregister mit Einzelzeichen
% weniger. `font=dense` ist dafür die falsche Antwort — es verkleinert die
% Schrift, während der Bedarf hier die Höhe ist.
%
% Als Faktoren (\newcommand) statt als Längen: \arraystretch ist ein
% Multiplikator auf die Zeilenhöhe, kein Mass. Deshalb auch bewusst NICHT
% über \ZSFDensityTables skaliert — der Bereichsfaktor wirkt auf die
% Abstände der Tabelle (\ZSFtableEdgePad, \extrarowheight), und ein Faktor
% mal einem Faktor ergäbe eine Stauchung, die niemand vorhersagt.
\newcommand{\ZSFtableArrayStretch}{1.04}
\newcommand{\ZSFtableArrayStretchRoomy}{1.28}
\newcommand{\ZSFtableArrayStretchTight}{0.92}

% Maximale Höhe für Bilder in figboxen (verhindert überhohe Aspect-Ratios).
\newlength{\ZSFfigMaxHeight}\setlength{\ZSFfigMaxHeight}{2.6cm}

% Höhe eines Bildes OHNE Box (\ZSFimage) in einer TABELLENZELLE. Eigener Wert
% und nicht \ZSFfigMaxHeight, weil die beiden verschiedene Grössenordnungen
% sind: eine Abbildung füllt einen Block, ein Zellenbild muss in eine
% Tabellenzeile passen.
%
% Der eigentliche Zweck ist die Gleichmässigkeit: Eine Spalte aus
% \ZSFimage{pfad}-Aufrufen ohne Höhenangabe ist damit von selbst einheitlich
% hoch. Stünde die Höhe stattdessen an jedem Aufruf, wäre dieselbe Entscheidung
% fünfmal geschrieben — und beim sechsten Bild einmal anders.
%
% Welcher der beiden Werte gilt, entscheidet NICHT das Bild, sondern der
% Baustein, in dem es steht: ZSFtable und valuegrid binden auf diesen Wert,
% figbox und splitbox auf \ZSFfigMaxHeight (Mechanik: \ZSF@imgBudget in
% styles/60_boxes.tex). Ein fester Default hier wäre an einem der beiden Orte
% still falsch — und still ist das Wort, das zählt: Ein zu klein gesetztes Bild
% bricht keinen Build und löst keine Warnung aus.
\newlength{\ZSFimageMaxHeight}\setlength{\ZSFimageMaxHeight}{1.1cm}

% Es gab hier ein \ZSFimagePadY — die senkrechte Luft, die ein Bild in einer
% Tabellenzelle braucht, weil ein \includegraphics Höhe und keine Tiefe hat
% und dadurch an der Oberkante der Zeile klebt. Der Wert ist ENTFALLEN: Es ist
% derselbe Satz wie »der Inhalt einer Zeile hält Abstand zur Zeilenkante«, und
% der gilt für einen \dfrac genauso. Zwei Register für eine Aussage sind eine
% Kopie (styles/README → Vereinigungstest); geblieben ist \ZSFtableCellPadY
% in 20_tables, gesetzt vom Regler `rows`.

% Nutzbarer Anteil der Zeilenbreite in mehrteiligen Bild-Layouts; der Rest ist
% der \hfill-Steg zwischen den Elementen (Bildreihe in \ZSFfig, Bild-neben-Text
% in \ZSFfigside) — EIN Wert, damit die beiden Layouts nicht um einen Punkt
% auseinanderdriften. Als Makro statt \newlength: es ist ein Faktor, kein
% Mass, und wird erst an der Verwendungsstelle gegen die dortige \linewidth
% ausgewertet.
\newcommand{\ZSFfigUsableFraction}{0.97}

% Abstand der TikZ-Spannklammer (\drawbrace, Opt-in 11_math_advanced) von der
% Grundlinie des annotierten Terms. Ohne ihn läuft die Klammer als waagrechte
% Linie mitten durch die Formel, weil die Marker auf der Grundlinie sitzen.
\newlength{\ZSFbraceDrop}\ZSFderive{\ZSFbraceDrop}{\ZSFspaceM}

%
%
%
%
%
%
%
%
%
%
%
%
%
%
%
\newlength{\ZSFmathRowSep}\ZSFderive{\ZSFmathRowSep}{\ZSFspaceSep}
\newlength{\ZSFmathColSep}\ZSFderive{\ZSFmathColSep}{\ZSFDensityText\ZSFspaceS}

% ── Der Abstand von der Spaltenoberkante zur ersten Grundlinie ───────
% \topskip beim Anfang einer Spalte, \splittopskip am Anfang einer
% fortgesetzten. Beide standen auf den Klassenvorgaben (9pt und 10pt) und
% waren damit die einzigen Masse, die JEDE Spalte anfasst und die niemand
% gesetzt hat — sie sind nicht »klein«, sie waren nur unsichtbar, weil eine
% Spalte meist mit einem Balken beginnt und der die Glue aufbraucht.
%
% EIN Register für beide: Es ist dieselbe Frage, nur einmal am Anfang und
% einmal in der Fortsetzung. Und bewusst an \baselineskip statt an der
% Dichte — der Wert muss die Oberlänge der ersten Zeile tragen, sonst rutscht
% sie nach oben aus dem Satzspiegel. Damit folgt er \ZSFLeadingFactor und
% nicht \ZSFDensityFactor, aus demselben Grund wie die Zeilenhöhe selbst.
% Gesetzt wird er in 70_document_settings, wo \baselineskip feststeht.
\newlength{\ZSFcolumnTopSkip}

% Tokens für die ruhige Box-Familie und ZSFlist.
% Der Einzug ist der einzige Wert, der sich zwischen nummerierter und
% unnummerierter Liste unterscheidet — eine Zahl mit Punkt braucht mehr Platz
% als ein Aufzählungszeichen. Deshalb zwei Werte unter EINEM Namen statt
% zweier Registersätze mit verschiedenen Präfixen.
\newlength{\ZSFstatementBarWidth}\setlength{\ZSFstatementBarWidth}{2pt}
\newlength{\ZSFlistMarkGap}         \ZSFderive{\ZSFlistMarkGap}{\ZSFspaceXS}
\newlength{\ZSFlistItemSep}         \ZSFderive{\ZSFlistItemSep}{\ZSFspaceXS}
\newlength{\ZSFlistLeftMargin}      \setlength{\ZSFlistLeftMargin}{1.2em}
\newlength{\ZSFlistLeftMarginOrdered}\setlength{\ZSFlistLeftMarginOrdered}{1.6em}
\newlength{\ZSFlistLabelSep}        \setlength{\ZSFlistLabelSep}{0.4em}

% Kleinstes vertikales Innenmass für Text an einer Kante: Zeilenluft in
% Tabellen (\extrarowheight in 70_document_settings) und die knappe Polsterung
% jeder Box mit `pady=tight` (leise Box, goalbox, valuegrid).
\newlength{\ZSFtextMinPad}\ZSFderive{\ZSFtextMinPad}{\ZSFDensityBoxes\ZSFspaceXS}

%
%
%
%
%
%
%
%
%
%
%
%
%
%
%
%
%
%
%
%
%
%
%
%
%
\newlength{\ZSFboxPadX}      \ZSFderive{\ZSFboxPadX}{\ZSFDensityBoxes\ZSFspaceS}
\newlength{\ZSFboxPadXBar}   \ZSFderive{\ZSFboxPadXBar}{\ZSFDensityBoxes\ZSFspaceM}
\newlength{\ZSFtableEdgePad} \ZSFderive{\ZSFtableEdgePad}{\ZSFDensityTables\ZSFspaceM}
\newlength{\ZSFtableEdgePadTight}
  \ZSFderive{\ZSFtableEdgePadTight}{\ZSFDensityTables\ZSFspaceS}

% Vertikal dasselbe: Innenabstand oben/unten und der Abstand NACH einer Box.
% Beide als Register, damit der Abstand nach einer Box nicht pro Box-Stil
% erneut ausgeschrieben wird.
\newlength{\ZSFboxPadY}      \ZSFderive{\ZSFboxPadY}{\ZSFDensityBoxes\ZSFspaceS}
\newlength{\ZSFBoxAfterSkip} \ZSFderive{\ZSFBoxAfterSkip}{\ZSFDensityBoxes\ZSFspaceS}

% Balken haben ihren eigenen vertikalen Rhythmus — bewusst getrennt vom
% Box-Innenabstand, damit sich Balkenhoehe und Boxpolsterung unabhaengig
% einstellen lassen.
\newlength{\ZSFbarPadY}      \ZSFderive{\ZSFbarPadY}{\ZSFDensityStructure\ZSFspaceS}
\newlength{\ZSFbarPadYTight} \ZSFderive{\ZSFbarPadYTight}{\ZSFDensityStructure\ZSFspaceXS}

% ------------------------------------------------------------
% Code-Längen — nur für die optionalen Code-Module benötigt
% (styles/65_code_style.tex + 67_code_comments.tex). Schadlos
% wenn die Module in preamble.tex deaktiviert sind.
% Threshold:  Schwelle für Auto-Overflow (Code+Kommentar > Threshold -> oben)
% Das Default-Register ist immutabel und dient \ResetCodeCommentThreshold
% als Anker; \SetCodeCommentThreshold modifiziert nur das "lebende" Register.
% ------------------------------------------------------------
% Zeilenhoehe im Code-Block. Eigenes Register statt einer nackten Zahl im
% Listings-Stil: Sie ist die einzige Zeilenhoehen-Entscheidung neben
% \ZSFLeadingFactor und soll wie dieser auffindbar und drehbar sein.
\providecommand{\ZSFcodeLeadingFactor}{0.8}
% \ZSFcodeMiniSkip ist ENTFALLEN. Es war der Abstand um ein Listing im
% globalen \lstset — der Stil CodeExpert setzte an derselben Stelle
% \ZSFspaceM, und welcher Wert galt, hing daran, ob der Autor den Stil
% angefordert hatte. Seit der Stil global gilt (65_code_style), gibt es nur
% noch eine Antwort, und sie kommt aus der Skala.
\newlength{\ZSFcodeFrameRule} \setlength{\ZSFcodeFrameRule}{0.4pt}
\newlength{\ZSFcodeXMargin}   \setlength{\ZSFcodeXMargin}{8pt}
\newlength{\ZSFcodeCommentThresholdDefault}
  \setlength{\ZSFcodeCommentThresholdDefault}{\maxdimen}
\newlength{\ZSFcodeCommentThreshold}\setlength{\ZSFcodeCommentThreshold}{\ZSFcodeCommentThresholdDefault}
\newlength{\ZSFcodeCommentGap}      \ZSFderive{\ZSFcodeCommentGap}{\ZSFspaceM}

% Valuegrid, Index, Readability und dokumentweite Overlays
% \ZSFvaluegridRowSep ist ENTFALLEN. Es war der Zeilenabstand des Wertegitters
% und damit ein zweites Register für dieselbe Frage, die der Regler `rows` an
% der ZSFtable beantwortet. Beide Tabellen-Bausteine lesen jetzt
% \ZSFtableCellPadY (20_tables) — ein Wert, zwei Engines.
\newlength{\ZSFreadabilityEmergencyStretch}\setlength{\ZSFreadabilityEmergencyStretch}{2.5em}
% Flatterzone im Fliesstext: schmaler als der ragged2e-Standard (0pt plus 2em),
% damit die schmalen 4 Spalten gut gefüllt bleiben. Wirkt nur über
% \ZSFReadableOn — Tabellen-Spalten (L/Y) und Index behalten den Standard.
\newlength{\ZSFreadabilityRaggedRightskip}\setlength{\ZSFreadabilityRaggedRightskip}{0pt plus 1em}
\newlength{\ZSFindexItemIndent}       \setlength{\ZSFindexItemIndent}{1.2em}
\newlength{\ZSFindexSubitemIndent}    \setlength{\ZSFindexSubitemIndent}{2em}
\newlength{\ZSFindexSubitemOffset}    \setlength{\ZSFindexSubitemOffset}{1em}
\newlength{\ZSFindexSubsubitemIndent} \setlength{\ZSFindexSubsubitemIndent}{2.6em}
\newlength{\ZSFindexSubsubitemOffset} \setlength{\ZSFindexSubsubitemOffset}{1.8em}
\newlength{\ZSFindexParStretch}       \setlength{\ZSFindexParStretch}{0.3pt}
% Deckkraft der Release-Kennung im Footer. Als Makro, weil \transparent einen
% einheitenlosen Faktor nimmt und ein solcher Wert sonst durch jede Pruefung faellt.
\providecommand{\ZSFfooterOpacity}{0.5}
\newlength{\ZSFfooterBottomInset}     \setlength{\ZSFfooterBottomInset}{1mm}
\newlength{\ZSFfooterRightInset}      \setlength{\ZSFfooterRightInset}{2mm}
% Freiraum am Satzspiegel-Fuss für das Footer-Overlay (Release-Kennung und
% Seitenzahl). Das Overlay zeichnet ausserhalb des Textflusses und reserviert
% deshalb KEINE Höhe — ohne diesen Freiraum überdruckt es die letzte Zeile
% einer voll gefüllten Spalte. Der Wert wird in 70_document_settings aus der
% Footer-Schrift gemessen und an geometry übergeben; hier steht nur das
% Register, damit alle Masse an einem Ort liegen.
\newlength{\ZSFfooterClearance}
\newlength{\ZSFitemizeLeftMargin}     \setlength{\ZSFitemizeLeftMargin}{3mm}
\newlength{\ZSFenumerateLeftMargin}   \setlength{\ZSFenumerateLeftMargin}{4mm}
\newlength{\ZSFidentityMarkerInset}   \setlength{\ZSFidentityMarkerInset}{2mm}

%
%
%
%
%
%
%
%
%
%
%
%
%
\newif\ifzsfAfterChapterBar
\zsfAfterChapterBarfalse
\newcommand{\ZSFMarkAfterChapterBar}{\global\zsfAfterChapterBartrue}
\newcommand{\ZSFConsumeAfterChapterBar}{\global\zsfAfterChapterBarfalse}

% Der Abstand VOR einem Block. Direkt nach einem Titelbalken der Balken-Abstand,
% sonst der normale Blockabstand.
%
% Gelesen wird die Balken-SPERRE (\ifzsfTitleBarPending, 60_boxes) und kein
% eigener Merker: Sie sagt bereits genau das, was hier gebraucht wird — »ein
% Balken steht da, und es kam noch kein Inhalt«. Ein zweiter Merker daneben
% wäre dieselbe Aussage ein zweites Mal, und die Sperre gilt für ALLE vier
% Balken, während der frühere \ifzsfAfterSubsectionBar nur die beiden
% Abschnittsebenen kannte — nach einem Kapitelbalken zog derselbe Block
% deshalb einen anderen Abstand.
\newcommand{\ZSFBoxBeforeSkip}{%
  \ifzsfTitleBarPending\ZSFbarAfterGap\else\ZSFspaceS\fi
}

% Kollabiert den before-skip der SubsectionBar auf \ZSFsubsectionAfterChapterGap,
% wenn sie direkt nach einer ChapterBar steht.
\newcommand{\ZSFCollapseAfterChapterBar}{%
  \ifzsfAfterChapterBar
    \vspace*{\dimexpr\ZSFsubsectionAfterChapterGap-\ZSFsubsectionBarBeforeSkip\relax}%
    \ZSFConsumeAfterChapterBar
  \fi
}

% ============================================================
% TEXT-BOX-BINDUNGEN
% Koppeln Fliesstext direkt an Boxen — verhindert visuellen Bruch.
% ============================================================

% Robustes Skip-/Penalty-Aufrollen (vmode-/hmode-aware)
\newcommand{\ZSFRobustUnskip}{%
  \ifvmode
    \ifdim\lastskip>0pt\relax\vskip-\lastskip\fi
    \ifnum\lastpenalty=0\else\unpenalty\fi
    \ifdim\lastskip>0pt\relax\vskip-\lastskip\fi
  \else
    \ifdim\lastskip>0pt\relax\unskip\fi
    \ifnum\lastpenalty=0\else\unpenalty\fi
    \ifdim\lastskip>0pt\relax\unskip\fi
  \fi
}

%
%
%
%
%
%
\makeatletter
% Hier deklariert, nicht in 60_boxes: Das Flag gehört zu diesen beiden Makros
% und muss vor ihnen existieren, egal in welcher Reihenfolge die Module
% wachsen. Gesetzt wird es von der formulabox.
\newif\ifZSF@inFormulaBox
%
%
%
%
%
%
%
%
%
%
%
%
%
\newcommand{\textVorBox}[1]{%
  \ifZSF@inFormulaBox
    \par
    \begingroup
    \setlength{\parskip}{0pt}%
    \noindent\raggedright\ZSFBoxNote{#1}\par
    \endgroup
    \vspace{\ZSFspaceXS}%
    \ZSFBoxAlign
  \else
    \par
    \addvspace{\ZSFspaceOuter}%
    \begingroup
    \setlength{\parskip}{0pt}%
    {\ZSFfontBoxBinding\color{TextVorBoxColor} \noindent #1\par}%
    \endgroup
    \nopagebreak
    \vspace{\ZSFspaceXS}%
    \vspace{-\ZSFspaceS}%
    \nopagebreak
  \fi
}

% Text der zur Box DARÜBER gehört (Ergänzung/Folgerung)
\newcommand{\textNachBox}[1]{%
  \ifZSF@inFormulaBox
    \par
    \vspace{\ZSFspaceXS}%
    \begingroup
    \setlength{\parskip}{0pt}%
    \noindent\raggedright\ZSFBoxNote{#1}\par
    \endgroup
    \ZSFBoxAlign
  \else
    \par
    \ZSFRobustUnskip
    \vspace{\ZSFspaceXS}%
    \begingroup
    \setlength{\parskip}{0pt}%
    \nopagebreak
    {\ZSFfontBoxBinding\color{TextNachBoxColor} \noindent #1\par}%
    \endgroup
    \addvspace{\ZSFspaceOuter}%
  \fi
}
\makeatother

%
%
%
%
%
%
\NewDocumentCommand{\ZSFgap}{O{M}}{\vspace{\csname ZSFspace#1\endcsname}}
% \ZSFSectionGap bleibt als eigener Name, weil er etwas anderes ausdrückt als
% eine Grösse: den Themenwechsel. Dass er aktuell derselbe Abstand wie
% \ZSFgap[L] ist, ist eine Einstellung dieser Zeile, keine Verdopplung.
\newcommand{\ZSFSectionGap}{\ZSFgap[L]}

% ------------------------------------------------------------
% Vertikaler Einschub INNERHALB einer Box
% ------------------------------------------------------------
% Formelnote, Trennlinie, Item-Heading, Folgerung und die Glieder einer
% Zielkette setzen alle denselben Einschub. Als Makro statt von Hand, damit der
% Abstand an der Dichte-Stellschraube hängt und der Rhythmus an einer Stelle
% änderbar bleibt.
%
% DER ABSTAND IST EIN OPTISCHER, KEIN LEIM. Das ist der ganze Punkt dieser drei
% Zeilen und der Grund, warum sie nicht `\vspace{#1}` heissen:
%
% Zwischen zwei Blöcken im vertikalen Satz liegt nicht nur die gesetzte Glue,
% sondern zusätzlich die Zeilenschaltung — TeX schiebt vor jede Box so viel
% Leim ein, dass der Grundlinienabstand \baselineskip ergibt. Der SICHTBARE
% Abstand ist damit `#1 + (\baselineskip - Höhe der folgenden Zeile)`. Solange
% alle Zeilen gleich hoch sind (Fliesstext), faellt das nicht auf. Sobald die
% Hoehen wechseln — Pillen einer Zielkette, Displaymathe, ein Bild —, wird der
% Rhythmus zum Nebenprodukt der Tinte: gemessen an der Zielkette der Showcase
% lagen zwischen benannt gleichen Elementen 4,1 bis 11,3 pt, obwohl der Autor
% genau zwei Werte gesetzt hatte.
%
% \nointerlineskip nimmt die Zeilenschaltung für den naechsten Block heraus.
% Danach ist der Abstand von der UNTERKANTE des vorigen Blocks zur OBERKANTE
% des naechsten genau #1 — unabhaengig davon, was in den beiden steht. Weil
% Hoehe und Tiefe einer Zeile der Tinte folgen (Oberlaenge, Unterlaenge), ist
% das zugleich der optische Abstand.
%
% Das `-\parskip` bleibt aus demselben Grund wie zuvor: Der folgende Block ist
% ein Absatz, und TeX setzt \parskip beim Absatzbeginn selbst. Ohne den Abzug
% käme er zu #1 hinzu.
%
% Die erste Zeile faengt den Fall ab, dass zwei Einschuebe direkt aufeinander
% treffen (\ZSFsep, unmittelbar gefolgt vom naechsten Kettenglied): Dann gilt
% der zweite, statt dass sich beide addieren.
\newcommand{\ZSFInterlude}[1]{%
  \par
  \ifdim\lastskip>0pt\relax \vskip-\lastskip \fi
  \vskip\dimexpr#1-\parskip\relax
  \nointerlineskip
}
% Trenner-Abstand um den Mittelpunkt im Index-Locator "6.2 · (S. 4)" (66_index).
% Bewusst als Makro statt \newlength: das em soll an der Verwendungsstelle in der
% Index-Schrift ausgewertet werden. Eine Länge würde es in der Präambel-Schrift
% einfrieren und den Abstand dadurch messbar verschieben.
\newcommand{\ZSFindexLocatorSep}{\hspace{0.2em}}
